% Section 5: the decoupled model, natural index, Toda structure (from archive/secS.tex, compressed to the statements used,
% with proof sketches; attributions from research/G1/RH/spectral/NOVELTY_TODA.md Sec. 5, second-hand (WebFetch), Rule C reading due).
\section{The natural index and the Toda structure}\label{sec:structure}

This section supplies three tools: the evenness of the integer layer in the natural index (which completes Theorem~\ref{thm:A}(ii)
and is the engine of Theorem~\ref{thm:C}), the Toda structure that organises the zeta layers as flows of the integer layers, and its
dispersionless limit, which gives the leading term of every zeta layer. In each case the mechanism is known in other settings and we
credit it where it is used; what is new is the setting, a fixed discrete measure accumulating at a hard edge, with no large parameter
other than $n$ and with zeta values as times.

\subsection{The decoupled model}\label{sec:decoupled}
The exponent $\alpha$ enters the jump data of Proposition~\ref{prop:N6c} in two ways: as the Bessel order $\nu$ in $a_k(\nu),b_k(\nu)$,
which are polynomials in $\nu^2$, and as the Szeg\H{o} exponent $\beta$ in $\rho^2=\phi^{2(\beta+1)}$ at the hard edge and in
$\phi^{-2(\beta+1)}$ at the soft edge, $\phi(z)=(z^{1/2}+(z-4c)^{1/2})/(2c^{1/2})$. Keeping them independent, the recursion and the extraction
of \S\ref{sec:recursion}, with the gauge of Lemma~\ref{lem:gauge}, define formal series
\begin{equation}\label{eq:decoupled}
  \frac{n^p\,b(n;\beta,\nu;\mu)}{c}=\sum_{\br}\mu^{\br}\,\cD_{\br}(n;\beta,\nu),\qquad
  B(n;\beta,\nu):=\cD_0(n;\beta,\nu)=\sum_{k\ge0}d_k(p,\beta,\nu)\,n^{-k},
\end{equation}
the \emph{decoupled model} $M(\beta,\nu)$ and its \emph{integer layer}; for $\beta=\nu=\alpha$ they are the layers of Theorem~\ref{thm:N6}. Each
$d_k$ is rational in $(p,\beta,\nu)$ and polynomial in $(\beta,\nu)$.

\subsection{Reflection, shift and evenness}\label{sec:shift-thm}
\begin{theorem}\label{thm:shift}
Identically in $(p,\beta,\nu)$, as formal series:
\begin{enumerate}
\item[(i)] \emph{(reflection)} $d_k(p,-\beta-2,\nu)=(-1)^kd_k(p,\beta,\nu)$ for all $k$;
\item[(ii)] \emph{(shift)} $B(n;\beta,\nu)=(n/N)^pF_p(N;\nu)$ with $N:=n+\frac{\beta+1}2$ and a formal series $F_p$ in $1/N$ that does not depend on
$\beta$; equivalently, for every real $\gamma$,
\begin{equation}\label{eq:FI}
  B(n+\gamma;\beta-2\gamma,\nu)=\Bigl(1+\frac\gamma n\Bigr)^pB(n;\beta,\nu);
\end{equation}
\item[(iii)] for $\br\in\{e_1,e_2,2e_1\}$ the series $(N/n)^{(1+|\br|_w)p}\cD_{\br}(n;\beta,\nu)$, re-expanded in $1/N$, does not depend on $\beta$.
\end{enumerate}
\end{theorem}

\begin{corollary}[evenness]\label{cor:even}
$F_p(N;\nu)=F_p(-N;\nu)$, and $F_p$ depends on $\nu$ only through $\nu^2$. For the lattice and every $p\ne1$ the integer layer of
Theorem~\ref{thm:N6} is, for every $K$, $\cD_{0,K}(n)=(n/N)^p[1+\sum_{2k\le K}f_{2k}(p,\alpha^2)N^{-2k}]+O(N^{-K-1})$ with $N=n+\frac\alpha2+\frac12$.
\end{corollary}

\begin{proof}
By~(i), $B(n;\beta,\nu)=B(-n;-\beta-2,\nu)$; the natural index of the right side is $-N$, so~(ii) gives $F_p(N)=F_p(-N)$. The $\nu$-dependence is
through $a_k(\nu),b_k(\nu)$ only.
\end{proof}

\begin{proof}[Proof of Theorem~\ref{thm:shift} (sketch)]
(i) Since $\sigma_2\mathsf T=\mathsf T\sigma_3$, $P_0^\infty$ and $P_1$ commute with $\sigma_2$. The core matrices $X_k$ and $\sigma_3\Xi_k\sigma_3$ lie in
$\operatorname{span}\{I,\sigma_2\}$ for even $k$ and in $\operatorname{span}\{\sigma_3,i\sigma_1\}$ for odd $k$, so conjugation by $\sigma_2$
multiplies them by $(-1)^k$, while $\sigma_2d^{\sigma_3}\sigma_2=d^{-\sigma_3}$. Hence $\sigma_2\Delta_k(\beta)n^{-k}\sigma_2=\Delta_k(-\beta-2)(-n)^{-k}$, the
recursion is covariant, and $(Q_1)_{12}(Q_1)_{21}$ is invariant. This grading is the device of Kuijlaars and Tibboel
\cite[Lemma~4.1]{KuijlaarsTibboel}, who use it to remove the odd powers of $1/n$ for varying weights; here it gives a covariance under
$(n,\beta)\mapsto(-n,-\beta-2)$.

(ii) No transformation of the formal recursion is known to relate $M(\beta,\nu)$ to itself (the Christoffel shift $x\sigma_{p,s}=\sigma_{p,s+p}$
moves $\beta$ and $\nu$ together). Instead we apply steepest descent to the degree-$n$ Riemann--Hilbert problem $Y$ of the lattice in the frame of a
non-integer index $m=n+\gamma$: put $U_m(z):=m^{pn\sigma_3}Y(zm^{-p})$ and $\tilde U:=CU_mz^{\gamma\sigma_3}C^{-1}$, $C=m^{p\gamma\sigma_3}$. The
residue at the node $(m/k)^p$ acquires the factor $k^{2p\gamma}$, so the band weight has exponent $\alpha-2\gamma$; the cut of $z^{\gamma\sigma_3}$
cancels the jump $e^{-2\pi im\sigma_3}$ on $(-\infty,0)$ because the degree $n$ is an integer; and in the hard-edge disc the lip entry is
$\omega^{(\alpha-2\gamma)}e^{\mp2\pi i\gamma}e^{-4\zeta^{1/2}}$, $\zeta=m^2f$, whose band limits are $e^{\mp i\pi\alpha}$. So the $m$-frame problem
is $M(\alpha-2\gamma,\alpha)$: Bessel order $\alpha$, Szeg\H{o} exponent $\alpha-2\gamma$. For $p>1$, $|\alpha|<1$ and $0\le\gamma<(1-|\alpha|)/2$ the
estimates of the Hurwitz analysis of \S\ref{sec:H-proof}, in which the index enters only as a real scale, transfer and give
$m^pb_n/c=\sum_{|\br|\le B}\mu_m^{\br}\cD_{\br,K}(m;\alpha-2\gamma,\alpha)+O(m^{-K-1}+m^{-2-(B+1)p}+m^{-1-(J+1)p})$. This transfer is written
out step by step in the notes of the repository; we do not repeat it. Now fix $p\in(1,2)\setminus\mathbb Q$, $\alpha\in(-\frac12,\frac12)$,
$\gamma\in(0,\frac14)$. Both frames express $n^pb_n/c$ as $\sum_ec_en^{-e}+o(n^{-E})$ with exponents $k+|\br|_wp$, pairwise distinct since $p$ is
irrational, so the coefficients agree; at integer exponents only $\br=0$ contributes, which is~\eqref{eq:FI} at $\beta=\alpha$. The coefficients of
the difference are rational in $(p,\alpha,\gamma)$ and vanish on an open set for infinitely many $p$, hence identically;
$\gamma=(\nu-\beta)/2$ gives (ii), and the equivalence with~\eqref{eq:FI} is algebra.

(iii) The same comparison at exponent $k+Jp$ gives $\sum_{|\br|_w=J}\hat\mu^{\br}[X_{\br}]_k=0$, $\hat\mu_j=2\zeta(-jp)/j$, with $X_{\br}$ the
difference of the two frames. For $J=1$, $\zeta(-p)\ne0$ on $(1,2)$; for $J=2$, $\zeta(-2p)[X_{e_2}]_k+4\zeta(-p)^2[X_{2e_1}]_k=0$ with brackets
rational in $p$, and $\zeta(-2p)/\zeta(-p)^2$ is not rational (it vanishes at the odd integers $p\ge3$). Weight $\ge3$ would need the linear
independence of the monomials $\prod_j\zeta(-jp)^{r_j}$ over $\mathbb Q(p)$, which we do not prove.
\end{proof}

Exact computer algebra confirms~\eqref{eq:FI} identically in $(p,\alpha,\beta,\gamma)$ to order $n^{-4}$, and $40$-digit evaluation to order
$n^{-8}$ at $32$ points including $p<1$; controls with a half shift, the wrong sign, or the Bessel order moved with $\beta$ fail (scripts
\texttt{m4\_shift\_exact}, \texttt{sigma\_frame\_exact}, \texttt{sigma\_frame\_num}). These checks confirm the algebraic consequence, not the
analysis. Evenness in a half-shifted index is known for varying weights with a string equation \cite[Theorem~5.2]{BleherIts},
\cite{EMP,KuijlaarsTibboel}, and in the semiclassical Laguerre case the index and the hard-edge exponent enter through $n+\alpha/2$ and
$\alpha^2$ \cite[Theorem~2.1]{FilipukVanAsscheZhang}. What is new here, as far as we know, is the decoupling of the Szeg\H{o} exponent from the
Bessel order, the identity~\eqref{eq:FI} for a fixed discrete measure, and its proof from a second frame at a non-integer index.

\subsection{The zeta layer as a Toda flow}\label{sec:toda}
In the unscaled variable the band factor $\exp(-\sum_k\mu_{n,k}z^k)$, $z=n^px$, is $\exp(\sum_kt_kx^k)$ with $n$-independent
\begin{equation}\label{eq:todatimes}
  t_k:=-\frac{2\zeta(-kp)}k ,
\end{equation}
the times of the Toda hierarchy. For a tilt $e^{\sum_jt_jx^j}\mu$ of a compactly supported measure, with Jacobi matrix $\mathcal J$ and
$h_n=\int\pi_n^2$, one has $\partial_{t_j}\log h_n=(\mathcal J^j)_{nn}$, $\partial_{t_j}b_n=\frac12b_n[(\mathcal J^j)_{nn}-(\mathcal J^j)_{n-1,n-1}]$ and
$\partial_{t_j}a_n=b_{n+1}(\mathcal J^j)_{n+1,n}-b_n(\mathcal J^j)_{n,n-1}$; this is classical (Moser, Flaschka; see \cite[Lemma~2.1,
Theorem~2.4]{VanAsscheToda}). Let $\cA_{\br}$ be the analogues of $\cD_{\br}$ for $n^pa_n/(2c)$, from the second moment of the same recursion and
the extraction of \cite{PaperII}, and $A:=\cA_0$. Substitute $\mu_j=-u_jn^{-jp}$ with $n$-independent $u_j$, put
$\hat b:=\sum_{\br}(-u)^{\br}n^{-|\br|_wp}\cD_{\br}(n)$, $\hat a$ likewise, and let $\hat{\mathcal J}$ be the tridiagonal matrix on $\mathbb Z$ with
$\hat{\mathcal J}_{ii}=2(1+\frac in)^{-p}\hat a(n+i)$, $\hat{\mathcal J}_{i,i-1}=\hat{\mathcal J}_{i-1,i}=\hat b_i:=(1+\frac in)^{-p}\hat b(n+i)$, the shift
$n\mapsto n+i$ acting at fixed $u$ and re-expanded in $1/n$.

\begin{theorem}[Toda structure]\label{thm:toda}
Let $p\ne1$, $1<s<1+2p$.
\begin{enumerate}
\item[(i)] For every $J$ and $\tau\in\R^J$, the tilted lattice $e^{\sum_j\tau_jx^j}\sigma_{p,s}$ satisfies Theorem~\ref{thm:N6}(i) with the same
layers and $\mu_{n,j}$ replaced by $\mu_{n,j}-\tau_jn^{-jp}$, and likewise for $n^pa_n/(2c)$ with the layers $\cA_{\br}$.
\item[(ii)] For every $j$, as identities of formal series,
$\partial_{u_j}\hat b=\tfrac12c^jn^{-jp}\hat b[(\hat{\mathcal J}^j)_{00}-(\hat{\mathcal J}^j)_{-1,-1}]$ and
$\partial_{u_j}(2\hat a)=c^jn^{-jp}[\hat b_1(\hat{\mathcal J}^j)_{10}-\hat b_0(\hat{\mathcal J}^j)_{0,-1}]$.
\item[(iii)] The identities of (ii) hold for $M(\beta,\nu)$ identically in $(p,\beta,\nu)$, and they determine every $\cD_{\br}$ and $\cA_{\br}$ from
the integer layers $B$ and $A$.
\end{enumerate}
\end{theorem}

\begin{proof}[Outline]
(i) The tilted lattice is $\sigma_{p,s,\ell}$ of \cite{PaperII} with $\ell(k)=\exp\sum_j\tau_jk^{-jp}$, analytic of every angle
\cite[Definition~1.1]{PaperII}; its band weight is that of the lattice times an entire polynomial tilt, which by Lemma~\ref{lem:gauge} changes
the band-factor parameters to $\mu_n-\tau n^{-p\cdot}$ and nothing else (constants depend on $\tau$; no uniformity is needed). (ii) Write
$u=\tau+t$ with $t$ as in~\eqref{eq:todatimes}. By (i) the tilted $b_m,a_m$ are polynomials in $u$ up to $O_u(m^{-p-L})$; by the Baire category
theorem these errors are uniformly bounded for large $m$ on some open ball, on which the Toda identities, integrated along $u_j$, give an identity
of polynomials in $(u,h)$ up to $O(n^{-p-L})$. Since $L$ is arbitrary, the coefficient of $h$ is (ii); no expansion is differentiated term by
term. (iii) The same argument for the Hurwitz lattices of Section~\ref{sec:hurwitz} gives (ii) on the open set $\{(\alpha+2\delta,\alpha)\}$, and the
coefficients are rational. The left side of (ii) at $u^{\gamma+e_j}$ is $(\gamma_j+1)$ times a new layer, the right side involves only lower
layers. The transfer of the estimates in (i) is written out in the notes of the repository.
\end{proof}

Exact computer algebra confirms (ii) on the layers of the recursion for $j\le3$ to second order in the times (scripts \texttt{toda\_layers},
\texttt{toda\_measures}); the variant with the shift at fixed scaled time $\mu$ fails at second order. Using the Toda equations with shifted
indices re-expanded in $1/n$ to generate the coefficients of an expansion is the method of \cite[Section~4]{EMP} for varying weights
$e^{-N(x^2/2+tx^{2\nu})}$; here the measure is fixed, the times enter as $u_jn^{-jp}$ with zeta base values, and the identities come from the
exact flows by a category argument rather than from a string equation. In this reading the zeta layer is not new data: it is a fixed,
$n$-independent Toda flow of the two integer layers, with times given by the zeta values.

\subsection{The index-shift law}\label{sec:indexshift}
Let $M(\beta,\nu)$ carry parameters $\mu_a=\hat\mu_an^{-pj_a}$ with exponents $j_1,\dots,j_r$, put
$\hat\varepsilon_a:=\frac{pj_a}2\binom{2j_a}{j_a}c^{j_a}\hat\mu_a$ ($\hat\varepsilon_j=\varepsilon_j$ for the lattice), and expand
$n^pb(n;\beta,\nu;\mu)/c=\sum_kn^{-k}b_k(\mu)$.

\begin{theorem}[index-shift law]\label{thm:indexshift}
Let $p\ne1$. Identically in $(p,\beta,\nu)$, $b_k$ is a polynomial of degree at most $k$ in $\mu$, and its part of degree exactly $k$ is the
$n^{-k}$ coefficient of $(n/x)^p$, where $x$ is the formal solution of $x=n-\frac1p\sum_a\hat\varepsilon_ax^{-pj_a}$. Equivalently, in natural
form the layer $\mu^{\br}$ is $(n/N)^p\hat\varepsilon^{\br}N^{-|\br|-S}[(p+1+S)_{|\br|-1}/(\br!\,p^{|\br|-1})+O(N^{-1})]$, $S:=p\sum_ar_aj_a$.
\end{theorem}

For one time this gives $\varepsilon_jN^{-1-jp}$, and $\frac{3p+1}{2p}$ for $\mu_1^2$ (Theorem~\ref{thm:N6}(ii)): at leading order in $1/n$ at each
order in the times, the zeta times act as an implicit shift of the index. \emph{Outline.} Grade the term $\hat\mu^{\br}n^{-k-S}$ by its
defect $k-|\br|$; the equations of Theorem~\ref{thm:toda}(ii) are in conservation form, so every term has defect $\ge0$, and the defect-zero parts
solve the dispersionless Toda hierarchy with initial data $a^0=2cn^{-p}$, $b^0=cn^{-p}$. These lie on the line $a=2b$ of a hard edge at $0$,
which is invariant, and on it every flow is a Hopf equation $\partial_{t_j}b=V_j(b)\partial_nb$, $V_j(b)=\frac j2\binom{2j}jb^j$, solved by
characteristics, $b=cx^{-p}$, $x=n+\sum_jt_jV_j(cx^{-p})$; Lagrange inversion gives the coefficients. This is the genus-zero mechanism of
random-matrix map enumeration (\cite[Section~2.3]{Ercolani}, \cite{EMP}; Lagrange inversion as in \cite[Section~26.4]{Bouttier}), and the
endpoints are Riemann invariants of the dispersionless Toda hierarchy (cf.~\cite{KonopelchenkoMAM}); the hard edge fixes one of them at $0$.
What differs here is the initial profile $b=cn^{-p}$ of a fixed measure, the reduction to the line $a=2b$, real and pole exponents, and
zeta-valued times; the mechanism itself is not new. The closed form agrees with every exact layer of the recursion to fourth order for
symbolic exponents and sixth order at rational parameters, and the Toda flow of the integer layers reproduces the dressings $g_j$ and $h_2$; a
``dressed'' index shift does not reproduce $h_2$, so the corrections to the leading term are dispersive (scripts \texttt{shift\_toda},
\texttt{shift\_exact}, \texttt{shift\_dressing}). The layer coefficients are rational in the exponents, so the identities continue to
non-integer exponents, in particular to the pole exponents of Theorem~\ref{thm:kappa}.

\begin{remark}[the diagonal layer and the parity of the zeta layers]\label{rem:diag}
Exact computer algebra on the second moment of the recursion gives, identically in $(p,\beta,\nu)$, $(M/n)^pA(n;\beta,\nu)=1+a_2M^{-2}+a_4M^{-4}+O(M^{-5})$
with $M=n+\frac\beta2+1=N+\frac12$, $a_2=-\frac{p+1}{24}+\frac{(p+1)(p+2)\nu^2}8$, no odd terms and no $\beta$ to this order (script
\texttt{toda\_layers}); evenness and $\beta$-independence beyond order four are supported only numerically (scripts \texttt{m4\_toda},
\texttt{m4\_diag}). So $b_n$ is centred at $N$ and $a_n$ at $N+\frac12$, as in \cite[Theorem~5.2]{BleherIts} for varying weights. If $A$ is even in
$M$, then $a_n-a_{n-1}$ is odd in $N$ and the second differences of $b_n^2$ are even, so by Theorem~\ref{thm:toda} the first-order zeta layers are
$\varepsilon_jN^{-1-jp}$ times even series and the $\varepsilon_1^2$ layer is $N^{-2-2p}$ times an even series, as found in
Theorem~\ref{thm:N6}(ii); the general monomial structure is open (Problem~\ref{prob:open}(b)). The factor $1/(2j-1)$ in $g_j$ matches the first
variation $\binom{2j-2}{j-1}/\binom{2j}j=j/(2(2j-1))$ of the dispersionless symbol $\frac1{2\pi}\int(a+2b\cos\theta)^j\diff\theta$ off the line $a=2b$
(heuristic).
\end{remark}
